% \clearpage
\section{Transformation Taxonomy}
\label{appendix:taxonomy}

\begin{table*}[htbp]
\centering
\caption{Taxonomy of Code Transformations and Their Cognitive Demands}
\label{tab:transformation_taxonomy}
\begin{adjustbox}{width=.75\linewidth,center}
\begin{tabularx}{\textwidth}{@{}l|l|X|X|X@{}}
\toprule
\textbf{Transformation Type} & \textbf{Level} & \textbf{Transformation Description} & \textbf{Reverse Task Goal} & \textbf{Cognitive Challenge} \\ \midrule
Identifiers Renaming\\(Variable, Class, Methods) & Lexical & Replaces meaningful identifiers with obfuscated names (e.g., \texttt{TotalNum} $\rightarrow$ \texttt{IILLLL}). & Generate readable, semantically appropriate identifiers (e.g., \texttt{Sum}, \texttt{x}). & Requires abstraction over identifier roles and naming conventions. \\ \addlinespace
Dead Code Insertion & Structural & Injects non-functional or misleading code blocks (e.g., \texttt{if (false) \{...\}}). & Remove irrelevant code while preserving original logic. & Requires distinguishing functional vs. non-functional code paths. \\ \addlinespace
Literals Encryption & Semantic & Encrypts literals (strings, integers) and embeds decryption logic within the code. & Replace encrypted literals with interpretable, semantically equivalent values. & Requires reasoning over encrypted values and the embedded decryption behavior. \\ \bottomrule
\end{tabularx}
\end{adjustbox}
\noindent\textit{Note:} For the \textbf{Identifier Renaming} transformation, all identifiers are replaced with visually similar patterns such as \texttt{IILLLL}, \texttt{IILL}, etc.
\end{table*}

\section{Prompt Templates}
\label{appendix:prompts}
This section provides the full templates for the prompts used in our experiments, including self-correction, reversibility assessment, and contrastive fine-tuning.

\subsection{Self-Correction Prompt Template}
When a model's generated output failed compilation, execution, or test cases, the following prompt (derived from our auto-correction script) was used to initiate self-correction.

\begin{tcolorbox}[colback=red!5, colframe=red!25, coltitle=black, title=Self-Correction Prompt, boxrule=0.4pt, left=2pt, right=2pt, top=2pt, bottom=2pt, fontupper=\small]
\textbf{System:} You are a code obfuscator. The previous obfuscated code did not work. Here is the code and the error. Please correct the code so it compiles and passes all tests.\\
\vspace{0.2cm}
\textbf{User:} \\
Original code (for reference):\\
\texttt{<Original Code>}\\
\vspace{0.1cm}
Previous obfuscated code:\\
\texttt{<Generated Code>}\\
\vspace{0.1cm}
Error message:\\
\texttt{<Compiler/Runtime Error or Failed Test Case>}\\
\vspace{0.1cm}
Please correct the code.
\end{tcolorbox}

\subsection{Deobfuscation Prompting Strategies}
This section provides the full text of the prompts used to assess the reversibility of fine-tuned models.

\subsubsection{Simple Prompt}~\\

\begin{tcolorbox}[colback=gray!5, colframe=gray!25, coltitle=black, boxrule=0.4pt, left=2pt, right=2pt, top=2pt, bottom=2pt, fontupper=\small]
\texttt{Deobfuscate the following Java code by <Obfuscation Technique> while preserving its functionality:}\\
\vspace{0.2cm}
\textbf{Input: \texttt{<Obfuscated code>}}
\end{tcolorbox}

\subsubsection{Few-Shot Prompt}~\\
\begin{tcolorbox}[colback=gray!5, colframe=gray!25, coltitle=black, boxrule=0.4pt, left=2pt, right=2pt, top=2pt, bottom=2pt, fontupper=\small]
\texttt{You are a code deobfuscator trained to deobfuscate Java code by renaming variables while preserving functionality... I will give you obfuscated code, you will analyze it, reverse engineer it with meaningful variable names, and provide human-readable code.}\\
\vspace{0.1cm}
\texttt{Example obfuscated code:}\\
\texttt{...[Example 1]...}\\
\vspace{0.1cm}
\texttt{Reverse engineered code:}\\
\texttt{...[Example 1 Deobfuscated]...}\\
\vspace{0.1cm}
\texttt{New obfuscated code:}\\
\texttt{<Obfuscated code>}
\end{tcolorbox}

\subsubsection{Chain-of-Thought (CoT) Prompt}~\\
\begin{tcolorbox}[colback=gray!5, colframe=gray!25, coltitle=black, boxrule=0.4pt, left=2pt, right=2pt, top=2pt, bottom=2pt, fontupper=\small]
\texttt{You are a Java code deobfuscator... reverse this process using a step-by-step approach.}\\
\texttt{IMPORTANT: You MUST follow each step completely...}\\
\texttt{Obfuscated Code to Deobfuscate: <Obfuscated code>}\\
\vspace{0.1cm}
\texttt{Step 1 - Analyze Program Purpose: ...}\\
\texttt{Step 2 - List ALL Variables and Their Roles: ...}\\
\texttt{Step 3 - Verify Complete Variable Mapping: ...}\\
\texttt{Step 4 - Apply Replacements to Generate Final Code...}
\end{tcolorbox}

\subsubsection{Task Description Without the Term ``Deobfuscation''}~\\
These prompts describe the inverse transformation explicitly and never use
the word ``deobfuscation''.

\begin{tcolorbox}[colback=gray!5, colframe=gray!25, coltitle=black,
  title=Identifier Renaming (explicit reverse), boxrule=0.4pt, left=2pt,
  right=2pt, top=2pt, bottom=2pt, fontupper=\small]
\texttt{Restore meaningful identifier names in the following Java code while
preserving its functionality:}\\
\vspace{0.2cm}
\textbf{Input: \texttt{<Obfuscated code>}}
\end{tcolorbox}

\begin{tcolorbox}[colback=gray!5, colframe=gray!25, coltitle=black,
  title=Dead Code Insertion (explicit reverse), boxrule=0.4pt, left=2pt,
  right=2pt, top=2pt, bottom=2pt, fontupper=\small]
\texttt{Remove non-functional or dead code blocks from the following Java
code and clean the program while preserving its functionality:}\\
\vspace{0.2cm}
\textbf{Input: \texttt{<Obfuscated code>}}
\end{tcolorbox}

\begin{tcolorbox}[colback=gray!5, colframe=gray!25, coltitle=black,
  title=Literal Encryption (explicit reverse), boxrule=0.4pt, left=2pt,
  right=2pt, top=2pt, bottom=2pt, fontupper=\small]
\texttt{Decrypt AES-encrypted literals in the following Java code and restore
their original values while preserving functionality:}\\
\vspace{0.2cm}
\textbf{Input: \texttt{<Obfuscated code>}}
\end{tcolorbox}

\subsection{Contrastive Fine-Tuning Prompt Templates}
The following three prompt formats were used to create the mixed dataset for Contrastive Fine-Tuning (CFT).

\subsubsection{Positive Classification}~\\
\begin{tcolorbox}[colback=green!5, colframe=green!25, coltitle=black, title=Positive Pair, boxrule=0.4pt, left=2pt, right=2pt, top=2pt, bottom=2pt, fontupper=\small]
\textbf{System:} You are an expert in code analysis and understanding semantic equivalence...\\
\textbf{User:} Analyze these two code snippets and determine if they are semantically equivalent...\\
Code Snippet 1: \texttt{\{original code\}}\\
Code Snippet 2: \texttt{\{Obfuscated code\}}\\
\textbf{Assistant:} Yes, these code snippets are semantically equivalent...
\end{tcolorbox}

\subsubsection{Negative Classification}~\\
\begin{tcolorbox}[colback=orange!5, colframe=orange!25, coltitle=black, title=Negative Pair, boxrule=0.4pt, left=2pt, right=2pt, top=2pt, bottom=2pt, fontupper=\small]
\textbf{System:} You are an expert in code analysis...\\
\textbf{User:} Analyze these two code snippets and determine if they are semantically equivalent...\\
Code Snippet 1: \texttt{\{original code\}}\\
Code Snippet 2: \texttt{\{Unrelated code\}}\\
\textbf{Assistant:} No, these code snippets are not semantically equivalent...
\end{tcolorbox}

\subsubsection{Forward Generation}~\\
\begin{tcolorbox}[colback=blue!5, colframe=blue!25, coltitle=black, title=Generation Task, boxrule=0.4pt, left=2pt, right=2pt, top=2pt, bottom=2pt, fontupper=\small]
\textbf{System:} You are an expert in code obfuscation techniques...\\
\textbf{User:} Obfuscate this Java code by Dead Code Insertion...: \texttt{\{Original Code\}}\\
\textbf{Assistant:} Here's the obfuscated version...: \texttt{\{Obfuscated code\}}...
\end{tcolorbox}

%\clearpage
\section{Detailed Experimental Results}
\label{appendix:results}

\subsection{CodeBLEU Scores for Deobfuscation Reversibility}

\begin{table*}[htbp]
\centering
\caption{Detailed CodeBLEU Scores for Deobfuscation Reversibility Test. The 'Baseline' is the similarity between original and tool-obfuscated code. The other columns show the similarity between the model's output and the obfuscated input. A successful deobfuscation would result in scores close to the baseline; instead, scores remain high, indicating failure.}
\label{tab:reversibility_heatmap_data}
\begin{tabular}{llccccc}
\toprule
\textbf{Model} & \textbf{Transformation} & \textbf{Baseline} & \textbf{Simple} & \textbf{Few-Shot} & \textbf{CoT} & \textbf{Augmented} \\
\midrule
GPT-4.1-Mini & Var. Renaming    & 0.50 & 0.7021 & 0.7124 & 0.7051 & 0.7628 \\
             & Literals Encryption & 0.22 & 0.7350 & 0.6800 & 0.6750 & 0.8000 \\
             & Dead Code Ins.   & 0.33 & 0.7500 & 0.6650 & 0.7550 & 0.7150 \\
\midrule
GPT-3.5-Turbo & Var. Renaming   & 0.50 & 0.7650 & 0.6500 & 0.7700 & 0.6950 \\
              & Literals Encryption & 0.22 & 0.8000 & 0.7450 & 0.6150 & 0.7300 \\
              & Dead Code Ins.  & 0.33 & 0.6850 & 0.7250 & 0.6900 & 0.7800 \\
\midrule
Mistral       & Var. Renaming   & 0.50 & 0.7500 & 0.6800 & 0.7550 & 0.7150 \\
              & Literals Encryption & 0.22 & 0.7850 & 0.7300 & 0.6300 & 0.7450 \\
              & Dead Code Ins.  & 0.33 & 0.7000 & 0.7100 & 0.7050 & 0.7650 \\
\midrule
\multicolumn{7}{l}{\textit{All other models showed similar patterns of high similarity to the obfuscated input.}} \\
\bottomrule
\end{tabular}
\end{table*}

\subsection{Error Distribution in Forward Task}
\label{appendix:error_analysis}

Of 2{,}705 failed samples (out of 5{,}400 total) after SFT and before
self-correction, 62\% produce incorrect outputs that compile and run but fail
test cases, 34\% crash at runtime, and 4\% fail to compile.
On Dead Code Insertion, incorrect outputs account for 69\% of failures.
On Literal Encryption, runtime crashes reach 45\% and incorrect outputs 52\%.
Table~\ref{tab:error_distribution} provides the full per-transformation
breakdown.

\begin{table}[htbp]
\centering
\caption{Distribution of error types for SFT forward failures before self-correction.}
\label{tab:error_distribution}
\resizebox{\columnwidth}{!}{%
\begin{tabular}{@{}l|ccc|c@{}}
\toprule
\textbf{Model} & \textbf{Compilation} & \textbf{Runtime} & \textbf{Incorrect Output} & \textbf{Total Failures} \\ \midrule
\multicolumn{5}{c}{\textit{Identifier Renaming}} \\ \midrule
\textit{Subtotal} & \textit{33 (6\%)} & \textit{118 (21\%)} & \textit{411 (73\%)} & \textit{562} \\ \midrule
\multicolumn{5}{c}{\textit{Dead Code Insertion}} \\ \midrule
\textit{Subtotal} & \textit{32 (4\%)} & \textit{230 (27\%)} & \textit{588 (69\%)} & \textit{850} \\ \midrule
\multicolumn{5}{c}{\textit{Literal Encryption}} \\ \midrule
\textit{Subtotal} & \textit{47 (4\%)} & \textit{578 (45\%)} & \textit{668 (52\%)} & \textit{1293} \\ \midrule
\textbf{Total} & \textbf{112 (4\%)} & \textbf{926 (34\%)} & \textbf{1667 (62\%)} & \textbf{2705} \\ \bottomrule
\end{tabular}%
}
\end{table}

\subsection{AST Structure Preservation Under Fine-Tuning}

\begin{table}[htbp]
\centering
\caption{Tree-sitter AST Analysis: Models Preserve Obfuscated Structure}
\label{tab:tree-sitter}
\small
\begin{tabular}{lcc}
\toprule
\textbf{Model} & \textbf{Identifier Retention ($J$)} & \textbf{Structural Similarity ($AST_{\text{struct}}$)} \\
\midrule
\textit{Baseline (Obf vs. Orig)} & 0.40 & 0.70 \\
\midrule
All Fine-Tuned Models & 0.999 -- 1.00 & 1.00 \\
\bottomrule
\end{tabular}
\end{table}

\section{Illustrative Example of Chain-of-Thought Failure}
\label{appendix:cot_failure}
One of the most striking findings was the models' inability to follow their own correct reasoning when prompted with a Chain-of-Thought (CoT) strategy. The model would correctly perform the intermediate reasoning steps, only to ignore its own conclusions and reproduce the obfuscated input verbatim in the final step. Figure~\ref{fig:cot_failure_example_appendix} provides a concrete example of this behavior.

\begin{figure*}[htbp]
\centering
\begin{tcolorbox}[colback=gray!5, colframe=gray!25, coltitle=black, title=Example of Cognitive Disconnect in Deobfuscation Task, boxrule=0.4pt, left=2pt, right=2pt, top=2pt, bottom=2pt]
\small\textbf{1. Obfuscated Code Input:}
\begin{lstlisting}[language=Java]
public class Calculator {
    public static void main(String[] args) {
        int l1l1l1 = 10;
        int l1l1l2 = 20;
        int l1l1l3 = l1l1l1 + l1l1l2;
        System.out.println("The sum is: " + l1l1l3);
    }
}
\end{lstlisting}
\hrule
\vspace{0.2cm}
\small\textbf{2. Model's Generated Reasoning (Correct Analysis):}
\begin{verbatim}
Step 1 - Analyze Program Purpose: This program initializes two integer variables, 
adds them together, and prints the result to the console.

Step 2 - List ALL Variables and Their Roles:
* l1l1l1: int - stores the first integer value (10) -> rename to "firstNumber"
* l1l1l2: int - stores the second integer value (20) -> rename to "secondNumber"
* l1l1l3: int - stores the sum of the first and second numbers -> rename to "sum"

Step 3 - Verify Complete Variable Mapping: All variables have been mapped.

Step 4 - Apply Replacements to Generate Final Code.
\end{verbatim}
\hrule
\vspace{0.2cm}
\small\textbf{3. Model's Final Output (Incorrect - Ignores Own Reasoning):}
\begin{lstlisting}[language=Java]
public class Calculator {
    public static void main(String[] args) {
        int l1l1l1 = 10;
        int l1l1l2 = 20;
        int l1l1l3 = l1l1l1 + l1l1l2;
        System.out.println("The sum is: " + l1l1l3);
    }
}
\end{lstlisting}
\end{tcolorbox}
\caption{A concrete example of a model (GPT-4.1-Mini) failing a deobfuscation task despite correctly executing the reasoning steps of a Chain-of-Thought prompt. The model correctly analyzes the code's purpose and maps variables but then fails to apply this reasoning, reproducing the obfuscated code verbatim. This demonstrates a fundamental disconnect between the model's reasoning and execution capabilities.}
\label{fig:cot_failure_example_appendix}
\end{figure*}

\section{Generalization Checks: Code Translation and Compilation}
\label{appendix:generalization}

To check whether the inversion curse extends beyond Java obfuscation, we
ran two smaller checks in different domains, summarized in
Table~\ref{tab:generalization_checks}. Both use a single API-based model
(GPT-4.1-mini), one fine-tuning run per direction, and temperature 0 for
every evaluation, matching the main study. Neither check runs the full BRF
protocol: there is no multi-model sweep, no four-prompting-strategy
comparison, and no chain-of-thought trace analysis. We report them as
corroborating evidence for the generality of the pattern, not as a third
application of BRF.

\textit{Python-to-Java translation.} We fine-tune GPT-4.1-mini on 1,000
Python-to-Java pairs from matched CodeNet Python800/Java250 problems, then
evaluate the untrained reverse direction (Java-to-Python) on 295--300
ConDefects/AtCoder programs using a $\geq$50\% test-pass threshold per
program.

\textit{C-to-assembly compilation.} We fine-tune GPT-4.1-mini on 1,000
C-to-assembly pairs from decompile-bench, then evaluate the untrained reverse
direction (assembly-to-C decompilation) on 164 HumanEval-Decompile (O0)
functions using re-execution against unit tests.

\begin{table}[htbp]
\centering
\caption{Generalization checks: reverse-direction accuracy before and after
fine-tuning on the forward direction only.}
\label{tab:generalization_checks}
\begin{adjustbox}{width=\columnwidth,center}
\begin{tabular}{@{}lcccc@{}}
\toprule
\textbf{Domain} & \textbf{Reverse Task} & \textbf{Baseline} & \textbf{After FT} & \textbf{Pass@3} \\ \midrule
Python$\to$Java & Java$\to$Python (test-pass) & 85.4\% & 63.3\% & 82.3\% \\
C$\to$asm & asm$\to$C (re-execution) & 16.5\% & 5.5\% & 8.5\% \\ \bottomrule
\end{tabular}
\end{adjustbox}
\end{table}

In both domains, fine-tuning on the forward direction only reduces reverse
accuracy below the unfinetuned baseline, the same pattern documented for Java
obfuscation in Section~\ref{sec:cognitive_specialization}. The two domains
diverge under additional generation attempts: Python-to-Java recovers most
of the gap at pass@3 (82.3\% against an 85.4\% baseline), while C-to-assembly
does not (8.5\% against a 16.5\% baseline). This is consistent with the
complexity-gating pattern reported for CFT in Section~\ref{contrastive}:
translation between two general-purpose languages has a comparatively
constrained, near-deterministic inverse, while decompiling assembly back to
C requires reconstructing a much larger and more ambiguous space of
equivalent source programs.
